\documentclass[12pt]{article}
\usepackage[utf8]{inputenc}
\usepackage[T1]{fontenc}
\usepackage[margin=1in]{geometry}
\usepackage{setspace}
\usepackage{xcolor}
\usepackage[colorlinks=true, urlcolor=blue!60!black, linkcolor=blue!60!black]{hyperref}
\setlength{\parindent}{2em}
\setlength{\parskip}{1.25em}
\renewcommand{\baselinestretch}{1.0}
\begin{document}
\begin{center}
{\LARGE \textit{Plessy v.\ Ferguson}, 163 U.S.\ 537 (1896)\par}
\vspace{0.8em}
{\large Case Brief --- Four-Page Edition\par}
\end{center}
\vspace{1em}

\section{Memory Jogger}

Over Justice Harlan's lone ``color-blind Constitution'' dissent, the Court upheld Louisiana's Separate Car Act under ``separate but equal,'' holding state-mandated segregation on intrastate railways a reasonable exercise of the police power consistent with the Thirteenth and Fourteenth Amendments (\href{https://supreme.justia.com/cases/federal/us/163/537/}{Justia}).

\section{Detailed Case Facts}

Louisiana's Separate Car Act of 1890 required railways to provide ``equal but separate accommodations for the white, and colored races,'' required conductors to assign passengers by race, and made a passenger who insisted on the wrong coach liable to a \$25 fine or 20 days in jail; only ``nurses attending children of the other race'' were exempt. The \textit{Comit\'e des Citoyens} of New Orleans organized a test case with the quiet cooperation of the East Louisiana Railway, which disliked the cost of separate cars. It recruited Homer Plessy, seven-eighths Caucasian and one-eighth African in ancestry, and hired a detective to ensure his arrest under the Act. On June 7, 1892, Plessy bought a first-class ticket from New Orleans to Covington, sat in the whites-only coach, refused to move, and was forcibly ejected, arrested, and jailed.

\section{Procedural History}

Before Judge John Howard Ferguson of the Orleans Parish criminal court, Plessy pleaded that the Act violated the Thirteenth and Fourteenth Amendments; the judge sustained the State's demurrer. The Louisiana Supreme Court upheld the Act and denied writs of prohibition and certiorari. Represented by Albion Tourg\'ee, Plessy took a writ of error to the U.S.\ Supreme Court, which affirmed on May 18, 1896. ``Separate but equal'' governed Jim Crow until \textit{Brown v.\ Board of Education} (1954).

\section{Judicial Votes}

\textbf{7--1}, affirmed.
\begin{itemize}
\item \textbf{Majority:} Justice Brown, joined by Chief Justice Fuller and Justices Field, Gray, Shiras, White, and Peckham.
\item \textbf{Dissent:} Justice Harlan.
\item \textbf{Not participating:} Justice Brewer.
\end{itemize}

\section{Holding}

State laws requiring ``equal but separate'' accommodations on intrastate conveyances violate neither the Thirteenth Amendment, which forbids only slavery and involuntary servitude, nor the Fourteenth, which secures political and civil equality but not ``social'' equality, so long as the separation is a ``reasonable'' exercise of the police power made with reference to ``the established usages, customs, and traditions of the people.'' Passages endorsing segregated schools and anti-miscegenation laws, and claiming ``[l]egislation is powerless to eradicate racial instincts,'' were dicta.

\section{Analysis of Opinions}

\textbf{Majority (Brown).} The Thirteenth Amendment claim was ``too clear for argument'': a law drawing ``merely a legal distinction'' between the races did not re-establish servitude. The Fourteenth Amendment sought legal equality but not to abolish distinctions of color or compel commingling. As proof, Brown cited segregated schools, including \textit{Roberts v.\ City of Boston} and Congress's own in the District of Columbia, and anti-miscegenation laws; he distinguished \textit{Strauder}'s jury rule as a political right. If Plessy was legally white he could sue for his reputation; if not, he had none. Regulations must be reasonable and not for ``the annoyance or oppression of a particular class,'' but legislatures have ``a large discretion.'' Any badge of inferiority arose ``solely'' from the construction Black citizens chose to put on the law.

\textbf{Harlan (dissenting).} A railroad is a ``public highway,'' so exclusion burdened civil rights and the liberty of locomotion. ``Every one knows'' the Act's purpose was to exclude Black passengers from white coaches. State-enforced separation imposed a ``badge of servitude,'' and the majority's reasonableness test would permit segregated streets, courtrooms, and juries. The decision would prove as pernicious as \textit{Dred Scott} and plant ``the seeds of race hate.'' Yet Harlan also affirmed that ``[t]he white race deems itself to be the dominant race'' and disparaged Chinese immigrants.

\section{Examples: Future Applications}

\begin{enumerate}
\item \textbf{Same-Side: Segregated Buses.} Separate front and rear bus sections with identical seats and fares fit \textit{Plessy}'s police-power rationale.
\item \textbf{Same-Side: Parks and Libraries.} Separate but comparable facilities were routinely sustained under \textit{Plessy}'s reasonableness test before 1954.
\item \textbf{Opposite-Side: Juries and Property.} Excluding Black citizens from juries or neighborhoods touches political and civil rights, invalid even under \textit{Plessy} (\textit{Strauder}; \textit{Buchanan v.\ Warley}).
\item \textbf{Opposite-Side: Unequal Law School.} A whites-only law school with no equal alternative fails ``separate \textit{but equal}'' itself (\textit{Gaines}; \textit{Sweatt v.\ Painter}).
\item \textbf{Fence-Sitter: Interstate Trains.} Dining-car curtains on interstate trains satisfy \textit{Plessy}'s equal protection deference but collide with the Dormant Commerce Clause (\textit{Morgan v.\ Virginia}, 1946).
\end{enumerate}

\section{Critique}

\textbf{Anti-canon.} \textit{Plessy} ranks with \textit{Dred Scott} in the anti-canon. Charles Black exposed the willful blindness of the ``badge of inferiority'' passage, since everyone knew segregation enforced white supremacy. Michael McConnell argues the Reconstruction framers understood the Amendment to forbid segregation, while Michael Klarman and Earl Maltz contend the ratifying public accepted Brown's civil--political--social distinction.

\textbf{Harlan's legacy.} Conservatives invoke the ``color-blind'' line against all race-conscious action; Goodwin Liu reads the dissent as anti-caste. Gabriel Chin's \textit{The Plessy Myth} warns against hagiography given Harlan's anti-Chinese rhetoric.

\section{Key Quotations}

\begin{quote}
``But in view of the constitution, in the eye of the law, there is in this country no superior, dominant, ruling class of citizens. There is no caste here. Our constitution is color-blind, and neither knows nor tolerates classes among citizens. In respect of civil rights, all citizens are equal before the law.''
\end{quote}
\noindent --- Justice Harlan, dissenting, 163 U.S.\ at 559.

\begin{quote}
``The object of the amendment was undoubtedly to enforce the absolute equality of the two races before the law, but, in the nature of things, it could not have been intended to abolish distinctions based upon color, or to enforce social, as distinguished from political, equality, or a commingling of the two races upon terms unsatisfactory to either.''
\end{quote}
\noindent --- Justice Brown, for the Court, 163 U.S.\ at 544.

\begin{quote}
``We consider the underlying fallacy of the plaintiff's argument to consist in the assumption that the enforced separation of the two races stamps the colored race with a badge of inferiority. If this be so, it is not by reason of anything found in the act, but solely because the colored race chooses to put that construction upon it.''
\end{quote}
\noindent --- Justice Brown, for the Court, 163 U.S.\ at 551.

\end{document}
